% =====================================================================
%  Báo cáo Assignment 06 — Phát triển các Hệ thống Thông minh
%  Nguyễn Duy Nghĩa · D23CTPM01 · B23DCCN600
%
%  Dựng bằng một lệnh:  python build.py
%  build.py chép dữ liệu hình, trích mã nguồn, sinh macro số liệu rồi gọi
%      latexmk -lualatex -shell-escape main.tex
%  Dùng LuaLaTeX (không phải XeLaTeX): bộ nhớ động, không vướng giới hạn
%  cố định của XeTeX khi pgfplots đọc các tệp dữ liệu lớn.
%  Mọi hình là TikZ/pgfplots, được externalize vào tikzcache/ nên chỉ biên
%  dịch lại khi hình đó đổi.
% =====================================================================
\documentclass[a4paper,13pt,oneside]{extreport}

% --------------------------- Ngôn ngữ & font -------------------------
\usepackage{fontspec}
\setmainfont{Times New Roman}
\setsansfont{Segoe UI}
% Consolas thiếu các ký hiệu toán (∘ ⊙ ∇) có trong docstring; luaotfload lấy chúng từ Cambria Math.
\directlua{luaotfload.add_fallback("toanhoc", {"Cambria Math:mode=node;"})}
\setmonofont{Consolas}[Scale=0.86, RawFeature={fallback=toanhoc}]
\usepackage{polyglossia}
\setmainlanguage{vietnamese}
% polyglossia đặt lại tên danh mục mỗi khi chọn ngôn ngữ, nên phải ghi đè trong \captionsvietnamese.
\gappto\captionsvietnamese{%
  \renewcommand{\listfigurename}{Danh mục Hình vẽ}%
  \renewcommand{\listtablename}{Danh mục Bảng biểu}%
  \renewcommand{\contentsname}{Mục lục}}

% --------------------------- Khổ trang -------------------------------
\usepackage[a4paper,left=3cm,right=2cm,top=2.5cm,bottom=2.5cm]{geometry}
\usepackage{setspace}
\onehalfspacing

% --------------------------- Toán học --------------------------------
\usepackage{amsmath,amssymb,amsfonts}
\usepackage{bm}

% --------------------------- Hình & bảng -----------------------------
\usepackage{graphicx}
\graphicspath{{assets/}}
\usepackage{booktabs,longtable,array,tabularx,multirow}
\usepackage{float}
\usepackage[table,dvipsnames]{xcolor}
\usepackage[font=small,labelfont=bf,labelsep=period]{caption}
\usepackage{subcaption}

% --------------------------- Danh sách -------------------------------
\usepackage{enumitem}
\setlist[itemize]{topsep=3pt,itemsep=2pt,parsep=0pt,leftmargin=1.5em}
\setlist[enumerate]{topsep=3pt,itemsep=2pt,parsep=0pt,leftmargin=1.7em}

% --------------------------- Màu & tiêu đề ---------------------------
\definecolor{ptitblue}{HTML}{123A75}
\definecolor{ptitred}{HTML}{C00000}
\definecolor{softgray}{HTML}{F2F4F7}

\usepackage{titlesec}
\titleformat{\chapter}[display]
  {\normalfont\sffamily\bfseries\color{ptitblue}}
  {\Large\chaptertitlename\ \thechapter}{6pt}{\LARGE}
\titlespacing*{\chapter}{0pt}{-18pt}{22pt}
\titleformat{\section}{\normalfont\large\sffamily\bfseries\color{ptitblue}}{\thesection}{0.6em}{}
\titleformat{\subsection}{\normalfont\sffamily\bfseries\color{ptitblue}}{\thesubsection}{0.6em}{}
\titleformat{\subsubsection}{\normalfont\bfseries\itshape}{\thesubsubsection}{0.6em}{}
\setcounter{secnumdepth}{3}
\setcounter{tocdepth}{2}

% --------------------------- Đầu trang / chân trang ------------------
\usepackage{fancyhdr}
\pagestyle{fancy}
\fancyhf{}
\renewcommand{\headrulewidth}{0.4pt}
\fancyhead[R]{\footnotesize\itshape Báo cáo Assignment 06: Mạng nơ-ron hồi quy trên dữ liệu chuỗi}
\fancyfoot[C]{\thepage}
\fancypagestyle{plain}{\fancyhf{}\fancyfoot[C]{\thepage}\renewcommand{\headrulewidth}{0pt}}

% --------------------------- Tài liệu tham khảo ----------------------
% biblatex không có vietnamese.lbx: để nguyên thì các chuỗi "in", "and", "pages" in ra thành tên khoá.
% Ánh xạ tiếng Việt sang bộ chuỗi tiếng Anh ("In:", "and", "pp.").
\usepackage[backend=biber,style=numeric,sorting=none,maxbibnames=6]{biblatex}
\DeclareLanguageMapping{vietnamese}{english}
\addbibresource{refs.bib}

% --------------------------- TikZ / pgfplots -------------------------
% =====================================================================
%  Style dùng chung cho mọi hình TikZ/pgfplots của báo cáo.
%  Mỗi hình nằm ở tikz/<tên>.tex và đọc dữ liệu từ figdata/*.dat.
% =====================================================================
\usepackage{tikz}
\usepackage{pgfplots}
\usepackage{pgfplotstable}
\pgfplotsset{compat=1.18}
\usetikzlibrary{external,positioning,arrows.meta,matrix,fit,calc,shapes.geometric,
                backgrounds,decorations.pathreplacing}
\usepgfplotslibrary{groupplots,colormaps}

% Externalize: mỗi hình biên dịch một lần thành tikzcache/<tên>.pdf.
\tikzexternalize[prefix=tikzcache/]
% Khi chạy bằng LuaLaTeX, external tự gọi lualatex cho từng hình, không cần khai system call.
% \hinh{tên}: nạp tikz/<tên>.tex, đặt tên tệp cache theo chính tên hình.
\newcommand{\hinh}[1]{\tikzsetnextfilename{#1}\input{tikz/#1}}

% Một màu cho mỗi mô hình, dùng thống nhất ở mọi hình (bảng màu an toàn cho người mù màu).
\definecolor{mZero}{HTML}{7F7F7F}   % Simple RNN
\definecolor{mOne}{HTML}{0072B2}    % LSTM
\definecolor{mTwo}{HTML}{E69F00}    % GRU
\definecolor{mThree}{HTML}{009E73}  % BiLSTM
\definecolor{accentRed}{HTML}{C0392B}
\definecolor{fwTorch}{HTML}{EE4C2C}   % PyTorch
\definecolor{fwKeras}{HTML}{0072B2}   % Keras
\definecolor{loyalC}{HTML}{009E73}    % KKBox: người ở lại
\definecolor{churnC}{HTML}{C0392B}    % KKBox: người rời bỏ
\definecolor{lightBlue}{HTML}{DCE8F5}
\definecolor{lightOrange}{HTML}{FCE9CF}
\definecolor{lightGreen}{HTML}{D5F0E6}
\definecolor{lightGray}{HTML}{EDEDED}

\pgfplotsset{
  a06/.style={
    width=7.2cm, height=5.2cm, grid=major, grid style={gray!25},
    tick label style={font=\footnotesize}, label style={font=\small},
    legend style={font=\footnotesize, draw=gray!50, fill=white, fill opacity=0.9, text opacity=1},
    title style={font=\small\bfseries},
    /pgf/number format/use comma, /pgf/number format/1000 sep={.}},
  m0/.style={mZero, thick, mark=none},
  m1/.style={mOne, thick, mark=none},
  m2/.style={mTwo, thick, mark=none},
  m3/.style={mThree, thick, mark=none},
  % Ảnh RGB từng điểm ảnh: tệp .dat có cột x y c, c dạng rgb255=r,g,b.
  anhRGB/.style={matrix plot*, mesh/cols=#1, mesh/color input=explicit,
                 point meta=explicit symbolic},
  anhRGB/.default=32,
  truc anh/.style={axis equal image, hide axis, enlargelimits=false, scale only axis},
}

\tikzset{
  khoi/.style={draw=ptitblue, fill=lightBlue, rounded corners=2pt, minimum height=7mm,
               align=center, font=\footnotesize, inner sep=3pt},
  khoiC/.style={khoi, fill=lightOrange, draw=mTwo!80!black},
  khoiG/.style={khoi, fill=lightGreen, draw=mThree!80!black},
  khoiX/.style={khoi, fill=lightGray, draw=gray},
  mui/.style={-{Stealth[length=2mm]}, thick, ptitblue},
  nhan/.style={font=\scriptsize, text=black!70},
  o/.style={draw=gray!70, minimum size=6mm, inner sep=0pt, font=\footnotesize},
}


% --------------------------- Mã nguồn --------------------------------
\usepackage[most]{tcolorbox}
\tcbuselibrary{external}
% tcolorbox vẽ khung bằng tikzpicture nội bộ; phải che khỏi externalize, nếu không job của từng hình sẽ lỗi.
\tcbset{shield externalize}
\usepackage{minted}
\setminted{fontsize=\footnotesize,linenos,numbersep=6pt,breaklines,frame=lines,
           framesep=2mm,baselinestretch=1.05,tabsize=4}
\usemintedstyle{friendly}
\renewcommand{\listingscaption}{Mã nguồn}
\renewcommand{\listoflistingscaption}{Danh mục Mã nguồn}
% \maNguon{module__ten}{chú thích}: chèn đoạn mã build.py đã trích từ src/Assignment 06/.
% Tên tệp chính là địa chỉ của đoạn mã: module a06/<module>.py (hoặc app/app.py), hàm hoặc lớp <ten>.
\newcommand{\maNguon}[2]{%
  \begin{listing}[H]
  \inputminted{python}{code/#1.py}
  \caption{#2}
  \label{code:#1}
  \end{listing}}
% \maNguonJS{ten}{chú thích}: hàm JavaScript <ten> trích từ app/static/app.js.
\newcommand{\maNguonJS}[2]{%
  \begin{listing}[H]
  \inputminted{javascript}{code/js__#1.js}
  \caption{#2}
  \label{code:js__#1}
  \end{listing}}

% --------------------------- Liên kết --------------------------------
\usepackage{hyperref}
\hypersetup{
  colorlinks=true, linkcolor=ptitblue, citecolor=ptitblue,
  urlcolor=ptitblue, filecolor=ptitblue,
  pdftitle={Assignment 06 — Mạng nơ-ron hồi quy trên dữ liệu chuỗi},
  pdfauthor={Nguyen Duy Nghia — B23DCCN600},
  pdfsubject={Phat trien cac He thong Thong minh},
  bookmarksnumbered=true
}
\usepackage{bookmark}
\usepackage{microtype}

% --------------------------- Lệnh riêng ------------------------------
\usepackage{etoolbox}
\newcommand{\nhanxet}[1]{\par\smallskip\noindent\textbf{Nhận xét:} #1\par\smallskip}
\newtcolorbox{phathien}[1]{
  breakable,
  colback=softgray, colframe=ptitblue, boxrule=0.7pt, arc=2pt,
  left=8pt, right=8pt, top=6pt, bottom=6pt,
  fonttitle=\bfseries\sffamily, title={#1}
}

% Số liệu: \SL{tệp/khoá/.../lá}, sinh tự động từ outputs/metrics/*.json bởi build.py.
% Thiếu khoá thì in ?? và ghi cảnh báo "Thieu so lieu"; build.py đọc main.log và báo lỗi, không nhận bản PDF đó.
% (Không dừng ngay tại chỗ vì mỗi hình externalize biên dịch lại cả tài liệu: một khoá thiếu ở chương khác
% sẽ làm hỏng mọi hình.)
\newcommand{\defSL}[2]{\expandafter\def\csname sl@#1\endcsname{#2}}
\newcommand{\SL}[1]{\ifcsname sl@#1\endcsname\csname sl@#1\endcsname\else
  \PackageWarning{solieu}{Thieu so lieu '#1'}\textbf{??}\fi}
\input{generated/so_lieu}


\renewcommand{\arraystretch}{1.16}
\newcolumntype{L}[1]{>{\raggedright\arraybackslash}p{#1}}
\newcolumntype{C}[1]{>{\centering\arraybackslash}p{#1}}
\newcolumntype{R}[1]{>{\raggedleft\arraybackslash}p{#1}}

\begin{document}

% =====================================================================
%  TRANG BÌA
% =====================================================================
% Bìa theo mẫu Assignment 04: khung ảnh lùi 6,5 mm khỏi mép giấy (197 × 284 mm), nội dung đặt tuyệt đối
% trong khung với cùng cỡ chữ và khoảng cách như bìa A04; chỉ đổi số bài và đề tài.
\begin{titlepage}
\thispagestyle{empty}
\tikzexternaldisable
\begin{tikzpicture}[remember picture,overlay]
  \node[anchor=north west, inner sep=0pt] at ([xshift=6.5mm,yshift=-6.5mm]current page.north west)
    {\includegraphics[width=197mm,height=284mm]{cover-frame.png}};
  \node[anchor=north, inner sep=0pt, text width=161mm, align=center]
    at ([yshift=-(6.5mm+20mm+6mm)]current page.north) {%
    {\bfseries\fontsize{15}{21}\selectfont HỌC VIỆN CÔNG NGHỆ BƯU CHÍNH VIỄN THÔNG\par}
    \vspace{4mm}
    {\bfseries\fontsize{14}{19}\selectfont KHOA CÔNG NGHỆ THÔNG TIN 1\par}
    \vspace{4mm}
    {\bfseries\fontsize{13.5}{18}\selectfont HỌC PHẦN PHÁT TRIỂN CÁC HỆ THỐNG THÔNG MINH\par}
    \vspace{12mm}
    \includegraphics[height=40mm]{ptit-logo.png}\par
    \vspace{10mm}
    {\bfseries\fontsize{22}{27}\selectfont\addfontfeature{LetterSpace=2}BÁO CÁO ASSIGNMENT 06\par}
    \vspace{5mm}
    {\bfseries\fontsize{13}{19.5}\selectfont ĐỀ TÀI: MẠNG NƠ-RON HỒI QUY RNN, LSTM, GRU VÀ BiLSTM\\
     TRÊN HAI TẬP DỮ LIỆU CHUỖI THỰC TẾ BẰNG PYTORCH VÀ KERAS\par}
    \vspace{13mm}
    {\fontsize{13}{16}\selectfont
     \renewcommand{\arraystretch}{1.8}
     \begin{tabular}{@{}p{52mm}l@{}}
       \textbf{Lớp:}                  & D23CTPM01 \\
       \textbf{Họ và tên:}            & Nguyễn Duy Nghĩa \\
       \textbf{Mã sinh viên:}         & B23DCCN600 \\
       \textbf{Giảng viên hướng dẫn:} & PGS.TS Trần Đình Quế \\
       \textbf{Học kỳ:}               & Học kỳ 1 năm học 2026 -- 2027 \\
     \end{tabular}\par}};
  \node[anchor=south, inner sep=0pt] at ([yshift=6.5mm+10mm+14mm]current page.south)
    {\bfseries\itshape\fontsize{13}{18}\selectfont Hà Nội, Tháng 09/2026};
\end{tikzpicture}
\tikzexternalenable
\null
\end{titlepage}

% =====================================================================
%  PHẦN ĐẦU
% =====================================================================
\pagenumbering{roman}
\setcounter{page}{1}
% =====================================================================
%  Lời mở đầu và bảng đối chiếu với sáu hạng mục của bài
% =====================================================================
\chapter*{Lời mở đầu}
\addcontentsline{toc}{chapter}{Lời mở đầu}
\markboth{Lời mở đầu}{Lời mở đầu}

Các bài trước làm việc với dữ liệu mà thứ tự không quan trọng: một bệnh nhân, một căn nhà, một bức ảnh. Assignment 06
chuyển sang dữ liệu chuỗi, nơi thứ tự chính là thông tin. Giá cổ phiếu hôm nay phụ thuộc hôm qua; một người sắp huỷ gói
nghe nhạc thường nghe thưa dần trước khi dừng hẳn. Mạng nơ-ron hồi quy được tạo ra cho đúng loại dữ liệu này: nó đọc
chuỗi từng bước và mang theo một trạng thái nhớ.

Bài dựng bốn kiến trúc hồi quy (Simple RNN, LSTM, GRU, BiLSTM) bằng cả PyTorch và Keras, huấn luyện trên hai tập dữ liệu
thật, tổng cộng 16 mô hình. Tập thứ nhất là giá cổ phiếu Amazon trong 27 năm, một bài hồi quy. Tập thứ hai là log nghe
nhạc theo ngày của \SL{kkbox/n_sample} người dùng KKBox, một bài phân loại người dùng rời bỏ dịch vụ. Chương 1 dựng lý
thuyết từ công thức, mỗi khái niệm là một hàm NumPy đối chiếu với PyTorch. Chương 2 phân tích và tiền xử lý hai tập.
Chương 3 và 4 huấn luyện 16 mô hình rồi so sánh hai framework. Chương 5 đóng gói các mô hình thành một Web App.

Hai tập cho hai kết quả trái ngược, và cả hai đều có ích. Trên KKBox, các mô hình học được một tín hiệu thật: người sắp
rời bỏ nghe nhiều hơn ở đầu tháng rồi giảm dần, và LSTM đạt ROC-AUC \SL{torch_kkbox/models/lstm/roc_auc} chỉ từ chuỗi
hành vi. Trên AMZN, không mô hình nào trong 16 mô hình thắng được mốc ngây thơ ``giá ngày mai bằng giá hôm nay''
(RMSE \SL{bench/naive/rmse} USD). Báo cáo trình bày kết quả này như nó là, kèm phân tích vì sao, thay vì dựa vào chỉ số
$R^2$ xấp xỉ 0,99 trên giá, một chỉ số trông rất đẹp nhưng mốc ngây thơ cũng đạt được.

Mọi con số trong báo cáo được sinh tự động từ tệp kết quả của notebook, mọi đoạn mã được trích thẳng từ mã nguồn, và mọi
hình được vẽ bằng TikZ từ dữ liệu, kể cả hai khung giao diện của Web App. Chạy lại notebook thì báo cáo đổi theo.
Phụ lục~\ref{ch:phuluc} ghi các lệnh để tái lập.

Em xin cảm ơn thầy \textbf{PGS.TS Trần Đình Quế} đã hướng dẫn học phần.

\vspace{8mm}
\begin{flushright}
\emph{Hà Nội, tháng 9 năm 2026}\par
\vspace{2mm}
\textbf{Nguyễn Duy Nghĩa}
\end{flushright}

\section*{Đối chiếu với các hạng mục của bài}

\begin{table}[H]
  \centering
  \small
  \begin{tabular}{L{4.4cm}L{1.3cm}L{4.2cm}L{4.2cm}}
    \toprule
    Hạng mục & Chương & Notebook & Mã nguồn \\
    \midrule
    1. Cơ sở lý thuyết RNN, kèm code & \ref{ch:lythuyet} & \texttt{00\_ly\_thuyet\_rnn} & \texttt{a06/theory.py} \\
    2. Phân tích và tiền xử lý hai tập dữ liệu chuỗi thực tế & \ref{ch:dulieu} & \texttt{01\_du\_lieu\_amzn}, \texttt{02\_du\_lieu\_kkbox} & \texttt{a06/data\_stock.py}, \texttt{a06/data\_churn.py} \\
    3. Hiện thực và huấn luyện bằng PyTorch & \ref{ch:pytorch} & \texttt{03\_pytorch\_amzn}, \texttt{04\_pytorch\_kkbox} & \texttt{a06/models\_torch.py}, \texttt{a06/train\_torch.py} \\
    4. Hiện thực Keras, đối chuẩn PyTorch và Keras & \ref{ch:keras} & \texttt{05\_keras}, \texttt{06\_doi\_chuan} & \texttt{a06/models\_keras.py}, \texttt{a06/bench.py} \\
    5. Đóng gói triển khai Web App & \ref{ch:webapp} & -- & \texttt{app/} \\
    6. Kết luận & \ref{ch:ketluan} & -- & -- \\
    \bottomrule
  \end{tabular}
\end{table}


\phantomsection
\addcontentsline{toc}{chapter}{Mục lục}
\tableofcontents
\cleardoublepage
\phantomsection
\addcontentsline{toc}{chapter}{Danh mục Hình vẽ}
\listoffigures
\cleardoublepage
\phantomsection
\addcontentsline{toc}{chapter}{Danh mục Bảng biểu}
\listoftables
\cleardoublepage
\phantomsection
\addcontentsline{toc}{chapter}{Danh mục Mã nguồn}
\listoflistings

% =====================================================================
%  PHẦN THÂN
% =====================================================================
\cleardoublepage
\pagenumbering{arabic}
\setcounter{page}{1}

% =====================================================================
%  Chương 1 — Hạng mục 1: cơ sở lý thuyết mạng nơ-ron hồi quy
%  Mã nguồn: a06/theory.py · Notebook: 00_ly_thuyet_rnn.ipynb
% =====================================================================
\chapter{Cơ sở lý thuyết mạng nơ-ron hồi quy}
\label{ch:lythuyet}

Chương này dựng lại các khái niệm của mạng hồi quy từ công thức. Mỗi khái niệm là một hàm NumPy trong
\texttt{a06/theory.py}; notebook 00 chạy hàm đó và so với lớp PyTorch tương ứng bằng \texttt{np.allclose}. Cả
\SL{theory/n_checks} phép đối chiếu đều cho kết quả khớp, nên các công thức dưới đây chính là những gì thư viện
đang tính.

\section{Dữ liệu chuỗi và tensor ba chiều}

Một mẫu dữ liệu chuỗi là một dãy vector $x_1, x_2, \dots, x_T$ với $x_t \in \mathbb{R}^D$. Gom $B$ mẫu lại được
tensor $X \in \mathbb{R}^{B \times T \times D}$ (Hình~\ref{fig:tensor}). Cả hai tập của bài đều có dạng này. AMZN
cho $X_{\text{train}}$ có shape $(\SL{theory/shapes/amzn/0}, \SL{theory/shapes/amzn/1}, \SL{theory/shapes/amzn/2})$:
mỗi mẫu là một cửa sổ 30 phiên giao dịch, mỗi phiên 8 đặc trưng. KKBox cho
$(\SL{theory/shapes/kkbox/0}, \SL{theory/shapes/kkbox/1}, \SL{theory/shapes/kkbox/2})$: mỗi mẫu là một người dùng,
31 ngày của tháng 3/2017, mỗi ngày 8 đặc trưng nghe nhạc.

\begin{figure}[H]
  \centering
  \hinh{f1_tensor}
  \caption{Tensor dữ liệu chuỗi; lát tô cam là một bước thời gian $x_t$}
  \label{fig:tensor}
\end{figure}

Điều làm dữ liệu chuỗi khác dữ liệu bảng là thứ tự. Giá hôm nay phụ thuộc giá hôm qua; một người sắp bỏ dịch vụ
thường nghe ít dần trong những ngày cuối. Giả định các mẫu độc lập và cùng phân phối (i.i.d.) vẫn đúng giữa các
mẫu, nhưng sai hoàn toàn giữa các bước trong cùng một mẫu. Mô hình cho chuỗi phải khai thác được sự phụ thuộc đó.

\section{Vì sao MLP và Conv1D chưa đủ}

Có hai cách dùng lại công cụ của các bài trước. Cách thứ nhất là làm phẳng chuỗi thành vector $T \cdot D$ chiều rồi
đưa vào MLP. Cách thứ hai là trượt một kernel Conv1D theo trục thời gian. Bảng~\ref{tab:mlp-rnn} đếm tham số của tầng
đầu tiên với $T = 30$, $D = 8$ và $H = 64$ đơn vị ẩn.

\begin{table}[H]
  \centering\small
  \caption{Tầng đầu của ba loại mô hình cho chuỗi $T = 30$, $D = 8$, $H = 64$}
  \label{tab:mlp-rnn}
  \begin{tabular}{lL{3.3cm}rL{3.2cm}L{3.2cm}}
    \toprule
    Mô hình & Số tham số & Giá trị & Độ dài chuỗi thay đổi & Nhớ quá khứ xa \\
    \midrule
    MLP (làm phẳng) & $(T \cdot D) H + H$ & \SL{theory/params/mlp} & không, $T$ cố định & mọi bước, nhưng không chia sẻ trọng số \\
    Conv1D, $K = 3$ & $K D H + H$ & \SL{theory/params/conv1d} & được & chỉ $K$ bước mỗi tầng \\
    Simple RNN & $H(D + H + 1)$ & \SL{theory/params/rnn} & được & qua trạng thái $h_t$ \\
    \bottomrule
  \end{tabular}
\end{table}

MLP học một trọng số riêng cho từng cặp (bước, đặc trưng). Nhân đôi độ dài chuỗi thì số tham số tầng đầu tăng từ
\SL{theory/params/mlp} lên \SL{theory/params/mlp_t60}, và điều học được ở bước 5 không dùng lại được ở bước 25.
Conv1D chia sẻ trọng số theo thời gian nhưng mỗi tầng chỉ nhìn $K$ bước, muốn thấy 30 bước phải chồng nhiều tầng.
RNN giải quyết cả hai: một bộ trọng số dùng cho mọi bước, và một trạng thái $h_t$ mang thông tin từ đầu chuỗi đến
hiện tại. Ba hàm đếm tham số nằm trong \texttt{theory.py}; notebook so từng công thức với \texttt{numel()} của
\texttt{nn.Linear}, \texttt{nn.Conv1d} và \texttt{nn.RNN}.

\maNguon{theory__count_params_rnn}{Số tham số của Simple RNN không phụ thuộc độ dài chuỗi}

\section{Simple RNN}

Elman~\cite{elman1990} đưa ra dạng đơn giản nhất: ở mỗi bước, trạng thái mới là hàm của đầu vào hiện tại và trạng
thái cũ,
\begin{equation}
  h_t = \tanh\left(W_{xh}\, x_t + W_{hh}\, h_{t-1} + b\right), \qquad h_0 = 0 .
  \label{eq:rnn}
\end{equation}
Mở vòng lặp theo thời gian thì RNN trở thành một mạng sâu $T$ tầng, nhưng mọi tầng dùng chung $W_{xh}$, $W_{hh}$ và
$b$ (Hình~\ref{fig:unroll}). Trạng thái cuối $h_T$ là bản tóm tắt cả chuỗi, và cả bốn mô hình của bài đều đưa
$h_T$ vào một đầu dự báo nhỏ.

\begin{figure}[H]
  \centering
  \hinh{f1_unroll}
  \caption{RNN ở dạng cuộn và dạng mở theo thời gian}
  \label{fig:unroll}
\end{figure}

\maNguon{theory__rnn_forward}{Lan truyền xuôi của Simple RNN theo công thức~\eqref{eq:rnn}}

Hàm trên khớp \texttt{nn.RNN} của PyTorch (kiểm tra \texttt{rnn\_forward}: \SL{theory/checks/rnn_forward}). PyTorch lưu
hai vector bias $b_{ih}$ và $b_{hh}$ nên tầng \texttt{nn.RNN(8, 64)} có \SL{theory/cell_params/rnn} tham số, nhiều hơn
công thức ở Bảng~\ref{tab:mlp-rnn} đúng $H = 64$. Khi so khớp, notebook cộng hai bias lại thành $b$.

\section{Lan truyền ngược theo thời gian và hai vấn đề của gradient}

Huấn luyện RNN dùng lan truyền ngược theo thời gian (BPTT)~\cite{werbos1990}: mở mạng ra $T$ bước rồi lan truyền
ngược như một mạng thường, cộng gradient của $W_{hh}$ từ mọi bước. Gọi $L$ là hàm mất mát tính từ $h_T$. Theo quy tắc
chuỗi,
\begin{equation}
  \frac{\partial L}{\partial h_t}
  = \frac{\partial L}{\partial h_T} \prod_{k=t+1}^{T} \frac{\partial h_k}{\partial h_{k-1}}
  = \frac{\partial L}{\partial h_T} \prod_{k=t+1}^{T} W_{hh}^\top\, \mathrm{diag}\!\left(1 - h_k^2\right).
  \label{eq:bptt}
\end{equation}
Gradient đi từ $h_T$ về $h_t$ bị nhân với $T - t$ ma trận Jacobi. Nếu chuẩn của các ma trận này nhỏ hơn 1, tích co lại
theo hàm mũ và gradient \emph{triệt tiêu}; nếu lớn hơn 1, tích \emph{bùng nổ}~\cite{bengio1994,pascanu2013}. Bỏ qua
đạo hàm của $\tanh$ (luôn không vượt 1) để thấy riêng vai trò của $W_{hh}$: với vector $v$ bất kỳ,
$\lVert (W_{hh}^\top)^t v \rVert$ tăng hay giảm theo bán kính phổ $\rho$ của $W_{hh}$.

\maNguon{theory__jacobian_product_norm}{Chuẩn của tích Jacobi với bán kính phổ $\rho$ cho trước}

Hình~\ref{fig:rho} chạy hàm này với ba giá trị $\rho$. Sau \SL{theory/rho_end/T} bước, $\rho = 0{,}9$ còn lại
\SL{theory/rho_end/r09}, còn $\rho = 1{,}1$ đã lên \SL{theory/rho_end/r11}. Chênh lệch chỉ $0{,}1$ ở bán kính phổ
thành chênh lệch hơn ba bậc độ lớn ở gradient.

\begin{figure}[H]
  \centering
  \hinh{f1_rho}
  \caption{Tích Jacobi co hoặc giãn theo hàm mũ tuỳ bán kính phổ của $W_{hh}$}
  \label{fig:rho}
\end{figure}

Thực tế có thêm $\tanh$ và trọng số khởi tạo ngẫu nhiên, nên notebook đo trực tiếp: dựng chuỗi ngẫu nhiên dài
\SL{theory/grad/T} bước, lấy $L = \sum h_T$, rồi dùng autograd đọc $\lVert\partial L/\partial h_t\rVert$ ở từng bước.
Để giữ được gradient của từng $h_t$, hàm mở vòng lặp bằng \texttt{RNNCell}/\texttt{LSTMCell} thay cho lớp RNN gộp.

\maNguon{theory__grad_norms_through_time}{Đo chuẩn gradient theo từng bước bằng autograd}

\begin{figure}[H]
  \centering
  \hinh{f1_gradnorm}
  \caption{Gradient của mất mát theo $h_t$, chuẩn hoá theo bước cuối (trục tung thang log)}
  \label{fig:gradnorm}
\end{figure}

Hình~\ref{fig:gradnorm} cho một kết quả đáng chú ý. Với Simple RNN, lùi 20 bước thì gradient chỉ còn
\SL{theory/grad/rnn_lag20} lần giá trị ở bước cuối; lùi 59 bước còn \SL{theory/grad/rnn_lag59}. LSTM với khởi tạo
mặc định của PyTorch cũng suy giảm gần như vậy: \SL{theory/grad/lstm_lag20} ở trễ 20. Lý do là bias cổng quên bằng 0
nên $f_t \approx 0{,}5$, và cell state bị nhân với khoảng $0{,}5$ ở mỗi bước. Khi đặt bias cổng quên bằng
\SL{theory/grad/forget_bias} như Jozefowicz và cộng sự đề xuất~\cite{jozefowicz2015}, $f_t$ tiến về gần 1 và gradient ở trễ
59 còn \SL{theory/grad/lstm3_lag59}, lớn hơn RNN hàng chục bậc. Vậy LSTM không tự động nhớ xa. Nó có một con đường để
gradient đi qua, và mô hình phải học mở con đường đó.

\section{Gradient clipping}

Triệt tiêu làm mô hình không học được phụ thuộc xa; bùng nổ thì nguy hiểm hơn, vì một bước cập nhật quá lớn có thể phá
hỏng trọng số đã học. Cách chữa thông dụng là cắt theo chuẩn~\cite{pascanu2013}: gộp mọi gradient thành một vector
$g$, nếu $\lVert g\rVert > c$ thì thay bằng $g \cdot c / \lVert g\rVert$. Độ dài bị chặn, còn hướng giữ nguyên
(Hình~\ref{fig:clip}). Cả 16 mô hình của bài dùng $c = 1$.

\maNguon{theory__clip_grad_norm}{Cắt gradient theo chuẩn, giữ hướng}

\begin{figure}[H]
  \centering
  \hinh{f1_clip}
  \caption{Gradient clipping chỉ đổi độ dài của $g$, không đổi hướng}
  \label{fig:clip}
\end{figure}

Trên hai ma trận gradient ngẫu nhiên có tổng chuẩn \SL{theory/clip/norm_before}, hàm trả về chuẩn
\SL{theory/clip/norm_after} và khớp \texttt{nn.utils.clip\_grad\_norm\_} (\SL{theory/checks/clip_grad_norm}).

\section{LSTM}

LSTM~\cite{hochreiter1997} thêm một trạng thái thứ hai là cell state $c_t$ và ba cổng điều khiển nó. Với
$z_t = [h_{t-1}; x_t]$,
\begin{align}
  f_t &= \sigma(W_f z_t + b_f), & i_t &= \sigma(W_i z_t + b_i), & o_t &= \sigma(W_o z_t + b_o), \nonumber\\
  \tilde c_t &= \tanh(W_c z_t + b_c), & c_t &= f_t \odot c_{t-1} + i_t \odot \tilde c_t, & h_t &= o_t \odot \tanh(c_t).
  \label{eq:lstm}
\end{align}
Cổng quên $f_t$ quyết định giữ bao nhiêu của $c_{t-1}$, cổng vào $i_t$ quyết định ghi bao nhiêu thông tin mới
$\tilde c_t$, cổng ra $o_t$ quyết định để lộ bao nhiêu ra $h_t$. Điểm mấu chốt là phép cộng ở dòng $c_t$
(Hình~\ref{fig:lstm}): $\partial c_t / \partial c_{t-1} = f_t$, không có ma trận trọng số và không có $\tanh$ trên
đường đi. Khi $f_t$ gần 1, gradient đi qua nhiều bước gần như nguyên vẹn.

\begin{figure}[H]
  \centering
  \hinh{f1_lstm}
  \caption{Tế bào LSTM: đường cộng của cell state ở trên, bốn phép biến đổi có cổng ở dưới}
  \label{fig:lstm}
\end{figure}

\maNguon{theory__lstm_cell}{Một bước LSTM theo~\eqref{eq:lstm}, thứ tự khối cổng như PyTorch}

\section{GRU}

GRU~\cite{cho2014} gộp cell state vào $h_t$ và chỉ dùng hai cổng: cổng reset $r_t$ quyết định dùng bao nhiêu trạng thái
cũ khi tạo ứng viên $n_t$, cổng cập nhật $z_t$ hoà trộn trạng thái cũ với ứng viên,
\begin{align}
  r_t &= \sigma(W_r z_t + b_r), \qquad z_t = \sigma(W_z z_t + b_z), \nonumber\\
  n_t &= \tanh\!\left(W_{in} x_t + r_t \odot (W_{hn} h_{t-1})\right), \qquad
  h_t = (1 - z_t) \odot n_t + z_t \odot h_{t-1}.
  \label{eq:gru}
\end{align}
Dòng cuối là một tổ hợp lồi: khi $z_t \to 1$ thì $h_t \approx h_{t-1}$, gradient đi thẳng qua giống đường cộng của
LSTM (Hình~\ref{fig:gru}).

\begin{figure}[H]
  \centering
  \hinh{f1_gru}
  \caption{Tế bào GRU với hai cổng}
  \label{fig:gru}
\end{figure}

\maNguon{theory__gru_cell}{Một bước GRU theo~\eqref{eq:gru}}

Bảng~\ref{tab:lstm-gru} so hai tế bào với cùng $D = 8$, $H = 64$. Số tham số đọc từ PyTorch.

\begin{table}[H]
  \centering\small
  \caption{So sánh Simple RNN, LSTM, GRU và BiLSTM ($D = 8$, $H = 64$, tham số của riêng tầng hồi quy trong PyTorch)}
  \label{tab:lstm-gru}
  \begin{tabular}{lcccr}
    \toprule
    Tế bào & Số khối trọng số & Trạng thái & Cổng & Tham số \\
    \midrule
    Simple RNN & 1 & $h_t$ & không & \SL{theory/cell_params/rnn} \\
    LSTM & 4 & $h_t$, $c_t$ & $f, i, o$ & \SL{theory/cell_params/lstm} \\
    GRU & 3 & $h_t$ & $r, z$ & \SL{theory/cell_params/gru} \\
    BiLSTM & $2 \times 4$ & $\overrightarrow{h}_t, \overleftarrow{h}_t$ & $f, i, o$ mỗi chiều & \SL{theory/cell_params/bilstm} \\
    \bottomrule
  \end{tabular}
\end{table}

\section{LSTM hai chiều}

BiLSTM~\cite{schuster1997} chạy hai LSTM độc lập, một đọc từ $t = 1$ đến $T$, một đọc từ $T$ về 1, rồi nối hai trạng
thái cuối: $[\overrightarrow{h}_T; \overleftarrow{h}_1] \in \mathbb{R}^{2H}$ (Hình~\ref{fig:bilstm}). Notebook kiểm tra
cách nối này trên \texttt{nn.LSTM(bidirectional=True)}: đầu ra có shape
$(\SL{theory/bilstm/out_shape/0}, \SL{theory/bilstm/out_shape/1}, \SL{theory/bilstm/out_shape/2})$, nửa đầu ở bước cuối
trùng $\overrightarrow{h}_T$ và nửa sau ở bước đầu trùng $\overleftarrow{h}_1$ (\SL{theory/checks/bilstm_concat}).

\begin{figure}[H]
  \centering
  \hinh{f1_bilstm}
  \caption{BiLSTM: vector tóm tắt là trạng thái cuối của hai chiều đọc}
  \label{fig:bilstm}
\end{figure}

Với bài dự báo giá, chiều ngược không làm lộ tương lai. Nó chỉ đọc lại 30 phiên của cửa sổ theo thứ tự ngược, mọi phiên
đều trước ngày cần dự báo. Chiều ngược thêm một góc nhìn khác lên cùng quá khứ, không thêm thông tin mới.

\section{Điều chuẩn cho mạng hồi quy}

Dropout thường đổi mặt nạ ở mỗi bước. Áp lên trạng thái $h_t$, nó xoá một phần trí nhớ khác nhau ở mỗi bước, và
thông tin khó đi xa. Gal và Ghahramani~\cite{gal2016} đề xuất dùng \emph{một} mặt nạ cho cả chuỗi (recurrent dropout):
đơn vị bị tắt thì tắt suốt chuỗi, đơn vị còn lại mang thông tin liền mạch (Hình~\ref{fig:dropout}, $p =
\SL{theory/mask_p}\,\%$).

\begin{figure}[H]
  \centering
  \hinh{f1_dropout}
  \caption{Mặt nạ dropout thường và recurrent dropout; ô đỏ là đơn vị bị tắt}
  \label{fig:dropout}
\end{figure}

\maNguon{theory__recurrent_dropout_masks}{Hai kiểu mặt nạ dropout cho chuỗi}

BatchNorm của các bài trước lấy thống kê theo lô và theo từng bước, không hợp với chuỗi có độ dài thay đổi.
LayerNorm~\cite{ba2016} chuẩn hoá theo chiều đặc trưng của \emph{từng mẫu}, nên dùng được ở mọi bước và mọi kích thước lô.
Hàm dưới khớp \texttt{nn.LayerNorm} (\SL{theory/checks/layer_norm}).

\maNguon{theory__layer_norm}{Layer Normalization theo từng mẫu}

Các mô hình ở Chương~\ref{ch:pytorch} chỉ có một tầng hồi quy 64 đơn vị và dừng sớm theo validation, nên bài không
dùng dropout hay LayerNorm. Hai kỹ thuật được trình bày ở đây vì chúng là công cụ chuẩn khi mô hình sâu hơn.

% =====================================================================
%  Chương 2 — Hạng mục 2: hai tập dữ liệu chuỗi thực tế
%  Mã nguồn: a06/data_stock.py, a06/data_churn.py · Notebook: 01_du_lieu_amzn, 02_du_lieu_kkbox
% =====================================================================
\chapter{Hai tập dữ liệu chuỗi thực tế}
\label{ch:dulieu}

Bài dùng hai tập dữ liệu thật, đủ lớn và thuộc hai loại bài toán khác nhau. AMZN là chuỗi giá cổ phiếu, một bài hồi
quy: dự báo giá đóng cửa của phiên kế tiếp. KKBox Churn là log nghe nhạc theo ngày của người dùng một dịch vụ trả phí,
một bài phân loại: người dùng này có bỏ dịch vụ hay không. Cả hai đều đưa về tensor $B \times T \times 8$ như
Chương~\ref{ch:lythuyet}, nên cùng một lớp mô hình chạy được trên cả hai.

\section{Giá cổ phiếu AMZN}

\subsection{Nguồn và quy mô}

Dữ liệu ngày của Amazon.com, Inc. (NASDAQ: AMZN) được tải từ Yahoo Finance~\cite{yahooamzn} bằng thư viện
\texttt{yfinance}, gồm \SL{amzn/n_raw} phiên từ \SL{amzn/start} (phiên đầu sau IPO) đến \SL{amzn/end}. Mỗi phiên có
giá mở cửa, cao nhất, thấp nhất, đóng cửa và khối lượng. Giá đã điều chỉnh theo các lần chia tách cổ phiếu, gần nhất là
lần chia 20:1 năm 2022. Tệp CSV được lưu trong kho mã nguồn để lần chạy sau dùng đúng dữ liệu này, không phụ thuộc
việc Yahoo cập nhật lại.

\maNguon{data_stock__add_features}{Thêm đặc trưng tài chính từ giá và khối lượng}

Tám đặc trưng đưa vào mô hình là Close, Open, High, Low, Volume, Return (log-return
$r_t = \ln(P_t/P_{t-1})$), MA20 (trung bình 20 phiên) và RSI14 (chỉ số sức mạnh tương đối của Wilder~\cite{wilder1978}).
MA50 và Volatility20 chỉ dùng để phân tích. Bỏ \SL{amzn/n_raw} $-$ \SL{amzn/n_feat} phiên đầu chưa đủ 50 phiên để tính
MA50, còn lại \SL{amzn/n_feat} phiên. Bảng~\ref{tab:amzn-desc} cho thống kê mô tả.

\begin{table}[H]
  \centering\small
  \caption{Thống kê mô tả của AMZN, \SL{amzn/n_feat} phiên}
  \label{tab:amzn-desc}
  \input{generated/tab_amzn_desc}
\end{table}

\begin{figure}[H]
  \centering
  \hinh{f2_price}
  \caption{Giá đóng cửa AMZN (thang log) với MA20, MA50 và khối lượng; vạch đứt là ranh giới train, val, test}
  \label{fig:price}
\end{figure}

Hình~\ref{fig:price} vẽ giá trên thang log. Trong 27 năm giá tăng hơn một nghìn lần, với hai lần sụt mạnh: bong bóng
dot-com 2000--2001 và đợt điều chỉnh 2022. Thang tuyến tính sẽ ép toàn bộ giai đoạn trước 2010 thành một đường nằm sát
trục hoành.

\subsection{Return, phân phối và biến động}

Hình~\ref{fig:return} xem xét log-return. Return trung bình \SL{amzn/ret/mean} mỗi phiên, độ lệch chuẩn
\SL{amzn/ret/std}. Phiên giảm mạnh nhất là \SL{amzn/ret/min_date} ($r = \SL{amzn/ret/min}$), tăng mạnh nhất là
\SL{amzn/ret/max_date} ($r = \SL{amzn/ret/max}$), đều trong giai đoạn dot-com. Phân phối có đuôi dày: độ nhọn dư
(excess kurtosis) là \SL{amzn/ret/kurt}, trong khi phân phối chuẩn có độ nhọn dư bằng 0. Tỉ lệ phiên lệch quá 4 độ lệch
chuẩn là \SL{amzn/ret/tail4}\,\%, nếu phân phối chuẩn thì chỉ \SL{amzn/ret/tail4_gauss}\,\%. Biến động
(ô c) tụ thành cụm: những giai đoạn dao động mạnh kéo dài nhiều tháng. Ô d cho thấy năm cột giá gần như trùng nhau
(tương quan gần 1). Chúng mang thông tin về mức giá, còn Return và RSI14 mang thông tin về thay đổi.

\begin{figure}[H]
  \centering
  \hinh{f2_return}
  \caption{Log-return, phân phối return so với phân phối chuẩn, biến động 20 phiên và tương quan giữa 8 đặc trưng}
  \label{fig:return}
\end{figure}

\subsection{Tính dừng: kiểm định ADF và tự tương quan}

Kiểm định Augmented Dickey--Fuller~\cite{dickey1979} có giả thuyết không là chuỗi có nghiệm đơn vị (không dừng).
Với giá đóng cửa, thống kê ADF là \SL{amzn/adf/close/stat}, $p = \SL{amzn/adf/close/p}$, không bác bỏ được giả thuyết:
giá không dừng. Với return, thống kê là \SL{amzn/adf/return/statSgn}, $p = \SL{amzn/adf/return/p}$: return dừng.
Hình~\ref{fig:acf} cho cùng kết luận. ACF của giá giảm rất chậm, ở trễ 40 vẫn còn \SL{amzn/acf_c40}. ACF của return
gần như nằm trong dải tin cậy $\pm\SL{amzn/acf_band}$; chỉ \SL{amzn/acf_r_out} trong 40 trễ vượt ra ngoài và đều rất
nhỏ.

\begin{figure}[H]
  \centering
  \hinh{f2_acf}
  \caption{ACF và PACF 40 trễ của giá và của return; dải xám là khoảng tin cậy 95\,\%}
  \label{fig:acf}
\end{figure}

Hai kết quả này quyết định cách đặt bài toán. Giá không dừng, nên mô hình học trên giá 1997--2016 sẽ gặp mức giá hoàn
toàn mới ở tập test: giá lớn nhất của test gấp \SL{amzn/p_max_ratio} lần giá lớn nhất của train. Một mô hình dự báo
thẳng mức giá đã chuẩn hoá Min-Max theo train sẽ phải ngoại suy ra ngoài khoảng $[0, 1]$ mà nó từng thấy. Return thì
dừng, nhưng gần như không có tự tương quan, tức là rất khó dự báo từ quá khứ của chính nó. Đó là điều giả thuyết thị
trường hiệu quả~\cite{fama1970} dự đoán, và Chương~\ref{ch:keras} sẽ thấy các mô hình va vào đúng giới hạn này.

\subsection{Cửa sổ trượt và chống rò rỉ}

Bài dự báo log-return phiên kế tiếp $r_{t+1}$ từ cửa sổ \SL{amzn/window} phiên, rồi quy đổi ra giá
$\hat P_{t+1} = P_t\, e^{\hat r_{t+1}}$ để báo sai số bằng USD. Trong mỗi cửa sổ, năm cột giá chia cho Close của ngày
cuối cửa sổ, nên mô hình chỉ thấy giá tương đối (ví dụ ``MA20 cao hơn giá hiện tại 3\,\%''), không thấy mức giá tuyệt đối.
Volume lấy log.

\maNguon{data_stock__make_windows}{Cửa sổ trượt với nhãn là return phiên kế tiếp}

Tổng cộng \SL{amzn/n_windows} cửa sổ, chia theo thời gian 70/15/15 và không xáo trộn (Hình~\ref{fig:window}). Tập
test là bốn năm cuối, từ \SL{amzn/split/test/start} đến \SL{amzn/split/test/end}. Mọi thống kê chuẩn hoá (trung bình, độ
lệch chuẩn của từng đặc trưng và của nhãn) chỉ tính trên phần train. Notebook 01 kiểm tra bằng \texttt{assert} rằng ngày
cuối của train đứng trước ngày đầu của val, và ngày cuối của val đứng trước ngày đầu của test.

\begin{figure}[H]
  \centering
  \hinh{f2_window}
  \caption{Cửa sổ trượt và cách chia theo thời gian của AMZN}
  \label{fig:window}
\end{figure}

\maNguon{data_stock__WindowScaler}{Chuẩn hoá đặc trưng và nhãn, fit chỉ trên train}

Độ lệch chuẩn của return ở ba phần cũng khác nhau: \SL{amzn/ret_std/train} ở train (có cả giai đoạn dot-com),
\SL{amzn/ret_std/val} ở val và \SL{amzn/ret_std/test} ở test. Mô hình học trên một thị trường biến động mạnh hơn nhiều so
với thị trường nó được đánh giá.

\section{KKBox Churn}

\subsection{Nguồn và định nghĩa nhãn}

KKBox là dịch vụ nghe nhạc trực tuyến lớn ở Đài Loan. Năm 2017, KKBox và hội nghị WSDM 2018 mở một cuộc thi trên
Kaggle~\cite{kkbox2017} với dữ liệu thật của người dùng đã ẩn danh. Bài dùng hai tệp:
\begin{itemize}
  \item \texttt{train\_v2.csv}: \SL{kkbox/n_labeled} người dùng, nhãn \texttt{is\_churn} bằng 1 nếu người dùng
        \emph{không} gia hạn gói thuê bao trong 30 ngày sau khi gói hết hạn vào tháng 3/2017;
  \item \texttt{user\_logs\_v2.csv} (1,4 GB, \SL{kkbox/log_rows} dòng): mỗi dòng là một ngày nghe của một người dùng
        trong tháng 3/2017, gồm số bài nghe dưới 25\,\%, 25--50\,\%, 50--75\,\%, 75--98,5\,\% và trên 98,5\,\% độ dài
        bài (\texttt{num\_25} đến \texttt{num\_100}), số bài khác nhau (\texttt{num\_unq}) và tổng số giây nghe
        (\texttt{total\_secs}).
\end{itemize}

\subsection{Lấy mẫu và dựng chuỗi}

Chỉ \SL{kkbox/n_with_log} người có cả nhãn lẫn ít nhất một ngày log. Từ đó bài lấy mẫu phân tầng
\SL{kkbox/n_sample} người (seed 42). Tệp log không được đọc trọn vào RAM: \texttt{pyarrow} đọc theo khối 64 MB, lượt
thứ nhất lấy tập người có log, lượt thứ hai giữ lại \SL{kkbox/log_rows_kept} dòng của người trong mẫu. Cả quá trình mất
khoảng một phút.

\maNguon{data_churn__sample_users}{Lấy mẫu phân tầng những người vừa có nhãn vừa có log}

Mỗi người dùng thành một ma trận $31 \times 8$: bảy cột log lấy $\log(1+x)$ vì phân phối đuôi rất dài, cột thứ tám là
cờ \texttt{active} (ngày đó có nghe hay không). Ngày không có log để 0. Một ngày có thể có nhiều dòng (nhiều thiết bị)
nên được cộng dồn. \SL{kkbox/secs_clipped} dòng có \texttt{total\_secs} âm hoặc vượt 86.400 giây (dài hơn một ngày) là
lỗi ghi log và được kẹp về $[0, 86.400]$.

\maNguon{data_churn__build_sequences}{Gom log theo ngày thành tensor $(N, 31, 8)$}

Chia 70/15/15 phân tầng theo nhãn: \SL{kkbox/n_train} người train, \SL{kkbox/n_val} val và \SL{kkbox/n_test} test.
Chuẩn hoá z-score từng đặc trưng, thống kê tính trên mọi ngày của mọi người trong train.

\subsection{Phân tích dữ liệu}

Nhãn lệch mạnh: \SL{kkbox/labels/churn} người churn trên \SL{kkbox/labels/total}, tức \SL{kkbox/churn_rate_sample}\,\%
(Hình~\ref{fig:kklabels}). Tỉ lệ này gần như giữ nguyên qua hai bước lọc (\SL{kkbox/churn_rate_labeled}\,\% trên toàn bộ
người có nhãn), nên mẫu đại diện được tập gốc. Một mô hình luôn đoán ``ở lại'' đã đạt accuracy
$100 - \SL{kkbox/churn_rate_sample} \,\%$, vì vậy accuracy không phải thước đo phù hợp; bài dùng F1, Recall, ROC-AUC và
PR-AUC.

\begin{figure}[H]
  \centering
  \hinh{f2_kkbox_labels}
  \caption{Phân bố nhãn trong mẫu và tỉ lệ churn qua các bước lọc}
  \label{fig:kklabels}
\end{figure}

Hình~\ref{fig:kktraj} là kết quả thú vị nhất của phần phân tích. Đầu tháng 3, nhóm sẽ churn nghe \emph{nhiều hơn} nhóm ở
lại ở cả ba đặc trưng. Trong tuần thứ hai (giữa ngày 6 và ngày 11) hai đường cắt nhau, rồi nhóm churn giảm đều đến cuối
tháng, còn nhóm ở lại gần như đi ngang. Tính trung bình cả
tháng, người churn có \SL{kkbox/active_days/churn} ngày nghe, người ở lại có \SL{kkbox/active_days/loyal} ngày. Tín hiệu
nằm ở \emph{xu hướng} theo thời gian chứ không nằm ở mức trung bình, và đó là loại tín hiệu mô hình hồi quy có thể đọc
được mà mô hình trên đặc trưng tổng hợp sẽ bỏ lỡ.

\begin{figure}[H]
  \centering
  \hinh{f2_kkbox_traj}
  \caption{Quỹ đạo trung bình 31 ngày của nhóm ở lại và nhóm churn}
  \label{fig:kktraj}
\end{figure}

\begin{figure}[H]
  \centering
  \begin{minipage}{0.47\textwidth}\centering\hinh{f2_kkbox_hist}\end{minipage}\hfill
  \begin{minipage}{0.5\textwidth}\centering\hinh{f2_kkbox_corr}\end{minipage}
  \caption{Trái: phân phối tổng giờ nghe trong tháng. Phải: tương quan giữa 8 đặc trưng}
  \label{fig:kkhist}
\end{figure}

Hình~\ref{fig:kkhist} (trái) cho thấy phân phối tổng giờ nghe của hai nhóm chồng lên nhau gần hết; trung vị là
\SL{kkbox/hours_med/churn} giờ với nhóm churn và \SL{kkbox/hours_med/loyal} giờ với nhóm ở lại. Các đặc trưng tương quan
mạnh với nhau (phải): \texttt{num\_100} và \texttt{total\_secs} có hệ số \SL{kkbox/corr_100_secs}. Bảng~\ref{tab:ttest}
so hai nhóm bằng kiểm định $t$ của Welch~\cite{welch1947}. Mọi khác biệt đều có ý nghĩa thống kê vì mẫu lớn, nhưng
độ chênh của trung bình nhỏ so với độ phân tán trong từng nhóm. Hành vi nghe chỉ phân biệt được một phần người churn;
phần còn lại phụ thuộc thông tin không có trong log, như phương thức thanh toán hay gói cước. Chương~\ref{ch:pytorch}
sẽ đo mức trần này.

\begin{table}[H]
  \centering\small
  \caption{Kiểm định $t$ Welch giữa nhóm ở lại và nhóm churn, trung bình 31 ngày của từng người}
  \label{tab:ttest}
  \input{generated/tab_ttest}
\end{table}

\maNguon{data_churn__stratified_split}{Chia train/val/test phân tầng theo nhãn}

% =====================================================================
%  Chương 3 — Hạng mục 3: hiện thực và huấn luyện bằng PyTorch
%  Mã nguồn: a06/models_torch.py, a06/train_torch.py, a06/bench.py
%  Notebook: 03_pytorch_amzn.ipynb, 04_pytorch_kkbox.ipynb
% =====================================================================
\chapter{Hiện thực và huấn luyện bằng PyTorch}
\label{ch:pytorch}

Chương này dựng bốn mô hình Simple RNN, LSTM, GRU và BiLSTM bằng PyTorch~\cite{paszke2019}, huấn luyện trên cả hai
tập của Chương~\ref{ch:dulieu} và đánh giá trên tập test. Bốn mô hình chỉ khác nhau ở tầng hồi quy; phần còn lại
(đầu dự báo, tối ưu, dữ liệu) giống hệt, để chênh lệch kết quả chỉ đến từ loại tế bào.

\section{Một lớp cho bốn mô hình}

Mô hình nhận tensor $(B, T, 8)$, cho qua tầng hồi quy 64 đơn vị, lấy vector tóm tắt $h_T$ (hoặc
$[\overrightarrow{h}_T; \overleftarrow{h}_1]$ với BiLSTM) rồi đưa qua đầu dự báo
$\mathrm{Linear}(\cdot, 32) \to \mathrm{ReLU} \to \mathrm{Linear}(32, 1)$. Đầu ra là một số: return đã chuẩn hoá với
AMZN, logit churn với KKBox. Cả bốn mô hình dùng chung một lớp; tham số \texttt{cell} chọn tầng hồi quy.

\maNguon{models_torch__PyTorchRNN}{Lớp chung cho Simple RNN, LSTM, GRU và BiLSTM}

\texttt{nn.LSTM} và \texttt{nn.GRU} gọi thẳng nhân cuDNN khi chạy trên GPU, nên vòng lặp theo thời gian không nằm
trong Python. Với BiLSTM, PyTorch trả về đầu ra hai chiều nối nhau ở mỗi bước; lớp trên cắt đúng
$\overrightarrow{h}_T$ ở bước cuối và $\overleftarrow{h}_1$ ở bước đầu, đúng như cách nối đã kiểm chứng ở
Chương~\ref{ch:lythuyet}.

\section{Vòng lặp huấn luyện}

Cấu hình giống nhau cho mọi mô hình: AdamW~\cite{loshchilov2019} với tốc độ học $10^{-3}$ và weight decay $10^{-4}$, cắt chuẩn gradient về 1 (Chương~\ref{ch:lythuyet}), dừng sớm khi val loss không
giảm sau một số epoch, và giữ lại trọng số của epoch có val loss nhỏ nhất. Bảng~\ref{tab:cfg} ghi phần khác nhau giữa hai
tập.

\begin{table}[H]
  \centering\small
  \caption{Cấu hình huấn luyện}
  \label{tab:cfg}
  \begin{tabular}{lll}
    \toprule
    & AMZN (hồi quy) & KKBox (phân loại) \\
    \midrule
    Hàm mất mát & MSE trên return chuẩn hoá & BCE trên logit, \texttt{pos\_weight} = số âm / số dương \\
    Lô, số epoch tối đa & 64, 60 & 512, 20 \\
    Dừng sớm & 8 epoch không cải thiện & 4 epoch không cải thiện \\
    Thiết bị & chọn bằng đo (Mục~\ref{sec:thietbi}) & GPU \\
    Seed & 42 & 42 \\
    \bottomrule
  \end{tabular}
\end{table}

Toàn bộ dữ liệu được chuyển lên thiết bị một lần dưới dạng tensor; mỗi epoch chỉ xáo chỉ số rồi cắt lô, không dùng
\texttt{DataLoader}. Với tập KKBox khoảng 100 MB, cách này bỏ được chi phí chép dữ liệu từ RAM sang GPU ở mỗi lô.

\maNguon{train_torch__fit}{Vòng lặp huấn luyện chung: AdamW, clip gradient, dừng sớm theo val loss}

Với KKBox, nhãn dương chỉ chiếm khoảng 9\,\%. Không cân bằng lại thì mô hình đạt loss thấp bằng cách đoán ``ở lại'' cho
tất cả. \texttt{pos\_weight} nhân mất mát của mẫu dương lên khoảng 10 lần, tương đương cân bằng hai lớp trong hàm mất mát.

\maNguon{train_torch__make_loss}{Hàm mất mát cho hai loại bài toán}

\subsection{Chọn thiết bị bằng đo}
\label{sec:thietbi}

Các bài trước cho thấy GPU không phải lúc nào cũng nhanh hơn với mô hình nhỏ và tập nhỏ. Vì vậy notebook 03 đo một
epoch LSTM trên cả hai thiết bị (sau một lượt khởi động) rồi chọn bên nhanh hơn. Kết quả: CPU
\SL{torch_amzn/probe/cpu} s/epoch, GPU \SL{torch_amzn/probe/cuda} s/epoch, nên AMZN cũng huấn luyện trên GPU. Nhân cuDNN
cho LSTM chạy cả 30 bước trong một lần gọi, còn CPU phải lặp từng bước.

\maNguon{train_torch__pick_device}{Chọn CPU hay GPU bằng cách đo một epoch}

\section{Kết quả trên AMZN}

\begin{figure}[H]
  \centering
  \hinh{f3_loss_amzn}
  \caption{Đường loss của bốn mô hình PyTorch trên AMZN (liền: train, đứt: val)}
  \label{fig:loss-amzn}
\end{figure}

Hình~\ref{fig:loss-amzn} cho thấy train loss giảm chậm nhưng đều, còn val loss gần như không đổi ngay từ những epoch
đầu. Val loss thấp hơn train loss nhiều vì nhãn được chuẩn hoá theo độ lệch chuẩn của train, trong khi return giai đoạn
val biến động ít hơn hẳn (Chương~\ref{ch:dulieu}); so hai đường với nhau thì không có ý nghĩa, chỉ xu hướng của từng
đường là có. BiLSTM có val loss nhỏ nhất ở epoch \SL{torch_amzn/models/bilstm/best_epoch}, tức là mọi epoch sau đó chỉ
làm mô hình khớp thêm nhiễu của tập train. Các mô hình còn lại đạt val loss nhỏ nhất ở epoch
\SL{torch_amzn/models/lstm/best_epoch} đến \SL{torch_amzn/models/rnn/best_epoch}.

\begin{table}[H]
  \centering\footnotesize
  \caption{Bốn mô hình PyTorch trên tập test AMZN (\SL{torch_amzn/test_start} đến \SL{torch_amzn/test_end})}
  \label{tab:torch-amzn}
  \input{generated/tab_torch_amzn}
\end{table}

Bảng~\ref{tab:torch-amzn} là kết quả trung thực nhất của bài, và nó không đẹp. RMSE của bốn mô hình nằm trong khoảng
\SL{torch_amzn/models/bilstm/rmse} đến \SL{torch_amzn/models/gru/rmse} USD, trong khi mốc ngây thơ ``giá ngày mai bằng
giá hôm nay'' đạt \SL{torch_amzn/naive/rmse} USD. Không mô hình nào thắng mốc ngây thơ. Cột $R^2$ trên giá cho
\SL{torch_amzn/models/lstm/r2} với mọi mô hình và với cả mốc ngây thơ, nên con số rất cao này không nói gì: giá hôm
nay đã giải thích gần hết phương sai của giá ngày mai. Cột $R^2$ trên return mới đo khả năng dự báo thật, và nó âm với
cả bốn mô hình. Tỉ lệ dự báo đúng chiều tăng hay giảm dao động quanh 50\,\%, từ
\SL{torch_amzn/models/rnn/dir_acc}\,\% đến \SL{torch_amzn/models/lstm/dir_acc}\,\%.

\begin{figure}[H]
  \centering
  \hinh{f3_pred_amzn}
  \caption{Giá thật, mốc ngây thơ và dự báo của bốn mô hình trên các phiên cuối của tập test}
  \label{fig:pred-amzn}
\end{figure}

Hình~\ref{fig:pred-amzn} cho thấy nhìn bằng mắt thì dự báo ``bám'' giá thật rất sát. Phóng to 30 phiên cuối (ô b) thì
thấy rõ lý do: đường dự báo gần như trùng mốc ngây thơ, tức là đường giá thật dịch sang phải một phiên. Mô hình học được
rằng return ngày mai có trung bình gần 0, và dự báo gần 0 là lựa chọn có sai số bình phương nhỏ nhất khi return không có
tự tương quan (Chương~\ref{ch:dulieu}). Chương~\ref{ch:keras} bàn tiếp giới hạn này.

\section{Kết quả trên KKBox}

\begin{figure}[H]
  \centering
  \hinh{f3_loss_kkbox}
  \caption{Đường loss của bốn mô hình PyTorch trên KKBox (liền: train, đứt: val)}
  \label{fig:loss-kkbox}
\end{figure}

KKBox là bài toán khác hẳn. Train loss và val loss giảm cùng nhau trong nhiều epoch (Hình~\ref{fig:loss-kkbox}); LSTM
đạt val loss tốt nhất ở epoch \SL{torch_kkbox/models/lstm/best_epoch}/\SL{torch_kkbox/models/lstm/epochs_run}. Mô hình
học được tín hiệu thật từ chuỗi hành vi.

Ngưỡng phân loại không cố định ở 0,5. Với mỗi mô hình, notebook chọn ngưỡng làm F1 lớn nhất trên \emph{validation},
rồi áp nguyên ngưỡng đó cho test. Chọn ngưỡng trên test sẽ làm kết quả test đẹp hơn thực tế.

\maNguon{bench__best_threshold}{Chọn ngưỡng theo F1 trên validation}

\begin{table}[H]
  \centering\footnotesize
  \caption{Bốn mô hình PyTorch trên tập test KKBox (\SL{kkbox/n_test} người dùng)}
  \label{tab:torch-kkbox}
  \input{generated/tab_torch_kkbox}
\end{table}

Ba mô hình có cổng (LSTM, GRU, BiLSTM) đạt ROC-AUC từ \SL{torch_kkbox/models/gru/roc_auc} đến
\SL{torch_kkbox/models/lstm/roc_auc}, cao hơn Simple RNN (\SL{torch_kkbox/models/rnn/roc_auc}). Khoảng cách không lớn
vì chuỗi chỉ dài 31 bước, chưa đủ dài để gradient của Simple RNN triệt tiêu hoàn toàn. LSTM có F1 cao nhất,
\SL{torch_kkbox/models/lstm/f1}\,\%, với Precision \SL{torch_kkbox/models/lstm/precision}\,\% và Recall
\SL{torch_kkbox/models/lstm/recall}\,\%: cứ ba người bị báo churn thì khoảng một người churn thật, và mô hình tìm ra
khoảng một phần ba số người churn.

\begin{figure}[H]
  \centering
  \hinh{f3_conf_kkbox}
  \caption{Ma trận nhầm lẫn trên test KKBox; màu ô theo tỉ lệ trong hàng}
  \label{fig:conf-kkbox}
\end{figure}

\begin{figure}[H]
  \centering
  \hinh{f3_roc_pr}
  \caption{Đường ROC và Precision--Recall trên test KKBox; vạch đứt là mô hình đoán ngẫu nhiên}
  \label{fig:roc-pr}
\end{figure}

Các con số này thấp so với bảng xếp hạng của cuộc thi, và có lý do. Các đội mạnh nhất dùng thêm lịch sử giao dịch (phương
thức thanh toán, tự động gia hạn, ngày hết hạn), là những đặc trưng dự báo churn trực tiếp. Bài này chỉ dùng log nghe
nhạc, vì trọng tâm của bài là mô hình hồi quy trên dữ liệu chuỗi. Hình~\ref{fig:roc-pr} cho thấy phần tín hiệu mà chuỗi hành vi
mang lại là có thật: đường PR nằm cao hơn nhiều so với mốc ngẫu nhiên \SL{torch_kkbox/pos_rate_test}\,\%, nhưng
ROC-AUC dừng ở khoảng 0,68. Phần tín hiệu còn lại không nằm trong việc người dùng nghe nhạc thế nào.

% =====================================================================
%  Chương 4 — Hạng mục 4: hiện thực Keras và đối chuẩn PyTorch với Keras
%  Mã nguồn: a06/models_keras.py, a06/bench.py · Notebook: 05_keras.ipynb, 06_doi_chuan.ipynb
% =====================================================================
\chapter{Keras và đối chuẩn hai framework}
\label{ch:keras}

Chương này viết lại bốn mô hình bằng Keras~\cite{chollet2015} (backend TensorFlow~\cite{abadi2016}), huấn luyện trên
cùng dữ liệu, rồi so 16 mô hình theo chất lượng, kích thước và tốc độ. Mọi chỉ số của cả hai framework được tính lại bằng
cùng một hàm trong \texttt{a06/bench.py}, từ dự báo đã lưu, để phép so không lệch vì cách tính.

\section{Bốn mô hình bằng Keras}

Kiến trúc giữ đúng như bản PyTorch: đầu vào \texttt{Input(T, 8)}, một tầng hồi quy 64 đơn vị (một trong
\texttt{SimpleRNN}, \texttt{LSTM}, \texttt{GRU} và \texttt{Bidirectional(LSTM)}), rồi hai tầng \texttt{Dense(32, relu)}
và \texttt{Dense(1)}. Tối ưu cũng giống:
AdamW với tốc độ học $10^{-3}$, weight decay $10^{-4}$, \texttt{clipnorm=1.0}.

\maNguon{models_keras__build_keras}{Bốn mô hình Keras dựng từ một hàm}

Có một khác biệt nhỏ về clipping. \texttt{clipnorm} của Keras cắt chuẩn của \emph{từng} tensor gradient, còn
\texttt{clip\_grad\_norm\_} của PyTorch cắt theo chuẩn gộp của mọi gradient. Với ngưỡng 1 và mô hình nhỏ như ở đây, cả hai
hiếm khi được kích hoạt sau vài epoch đầu, nên bài giữ mặc định của mỗi framework.

Keras không có tham số \texttt{pos\_weight}. Cách tương đương là trọng số mẫu: mẫu dương nhận trọng số bằng số âm chia số
dương, mẫu âm nhận 1. Trọng số được truyền cho cả tập train lẫn tập validation, để val loss mà \texttt{EarlyStopping}
theo dõi đo đúng đại lượng như val loss của PyTorch.

\maNguon{models_keras__sample_weights}{Trọng số mẫu thay cho \texttt{pos\_weight}}

\maNguon{models_keras__fit_keras}{Huấn luyện một mô hình Keras với cùng lô và cùng dừng sớm}

\subsection{Chênh lệch số tham số}

Cùng kiến trúc nhưng số tham số không bằng nhau (Hình~\ref{fig:size}a). PyTorch lưu hai vector bias cho mỗi khối cổng
($b_{ih}$ và $b_{hh}$), Keras chỉ lưu một. Vì vậy Keras ít hơn đúng $H = 64$ tham số với Simple RNN
(\SL{bench/cmp/amzn_rnn/param_diff}), $4H$ với LSTM (\SL{bench/cmp/amzn_lstm/param_diff}) và $8H$ với BiLSTM
(\SL{bench/cmp/amzn_bilstm/param_diff}). GRU là ngoại lệ: Keras mặc định \texttt{reset\_after=True}, áp cổng reset sau
phép nhân ma trận giống PyTorch, nên cần hai vector bias và số tham số bằng nhau
(\SL{bench/models/amzn_torch_gru/params}). Bài giữ nguyên chênh lệch này thay vì ép bằng, vì nó không đổi khả năng biểu
diễn: tổng hai bias chỉ là một bias.

\begin{figure}[H]
  \centering
  \hinh{f4_size}
  \caption{Số tham số và kích thước tệp mô hình của hai framework}
  \label{fig:size}
\end{figure}

Kích thước tệp chênh nhiều hơn số tham số. Tệp \texttt{.pth} chỉ chứa \texttt{state\_dict}, còn tệp \texttt{.keras}
mặc định lưu thêm cấu hình mô hình và trạng thái của optimizer (hai moment của AdamW cho mỗi trọng số), nên nặng gấp
khoảng ba lần: \SL{bench/models/amzn_keras_lstm/size_kb} KB so với \SL{bench/models/amzn_torch_lstm/size_kb} KB với LSTM.
Nhờ vậy tệp Keras huấn luyện tiếp được ngay, nhưng khi chỉ cần phục vụ thì phần optimizer là thừa.

\section{Huấn luyện song song trên hai thiết bị}

TensorFlow từ bản 2.11 không còn hỗ trợ GPU trên Windows, nên Keras chạy trên CPU (16 luồng) trong một tiến trình riêng,
\emph{cùng lúc} với notebook PyTorch dùng GPU. Hai bên không phải chờ nhau, và CPU 32 luồng của máy được chia đôi cho hai
tiến trình. Hình~\ref{fig:loss-keras-amzn} và~\ref{fig:loss-keras-kkbox} là đường loss của bốn mô hình Keras; hình dạng
giống bản PyTorch ở Chương~\ref{ch:pytorch}: trên AMZN val loss đứng yên ngay từ đầu, trên KKBox cả hai đường cùng giảm.

\begin{figure}[H]
  \centering
  \hinh{f4_loss_keras_amzn}
  \caption{Đường loss của bốn mô hình Keras trên AMZN (liền: train, đứt: val)}
  \label{fig:loss-keras-amzn}
\end{figure}

\begin{figure}[H]
  \centering
  \hinh{f4_loss_keras_kkbox}
  \caption{Đường loss của bốn mô hình Keras trên KKBox (liền: train, đứt: val)}
  \label{fig:loss-keras-kkbox}
\end{figure}

Thời gian huấn luyện phản ánh chủ yếu thiết bị, không phản ánh framework. Trên KKBox, một epoch BiLSTM mất
\SL{bench/models/kkbox_keras_bilstm/epoch_s} giây với Keras trên CPU và \SL{bench/models/kkbox_torch_bilstm/epoch_s} giây
với PyTorch trên GPU. Vì vậy phần so tốc độ bên dưới đo suy luận trên cùng một thiết bị.

\section{Đo độ trễ suy luận}

Để so công bằng, độ trễ của cả 16 mô hình được đo trên CPU với đúng một luồng. Mỗi mô hình chạy khởi động 20 lần, rồi đo
200 lần với lô một mẫu và lấy trung vị (độ trễ), và đo 30 lần với lô 256 mẫu (thông lượng). Mô hình Keras chạy qua
\texttt{tf.function}, tức đồ thị đã biên dịch, là cách dùng khi phục vụ thật; gọi \texttt{model(x)} ở chế độ eager chậm
hơn khoảng một trăm lần vì vòng lặp theo thời gian chạy từng phép tính một.

\maNguon{bench___median_time}{Đo thời gian: khởi động rồi lấy trung vị}
\maNguon{bench__latency}{Độ trễ với lô 1 và thông lượng với lô 256}

Khi viết notebook này, lần đo đầu tiên cho PyTorch chậm bất thường, 1,5 đến 5 ms mỗi mẫu. Nguyên nhân là mô hình PyTorch và
Keras được đo xen kẽ nhau: sau mỗi lần gọi, các luồng của TensorFlow tiếp tục quay chờ việc một lúc và giành mất lõi CPU
duy nhất của PyTorch. Đo hết PyTorch rồi mới đo Keras thì kết quả ổn định. Đây là một ví dụ nhỏ cho thấy đối chuẩn hiệu
năng dễ sai như thế nào nếu không kiểm soát môi trường.

\section{Bảng đối chuẩn 16 mô hình}

\begin{table}[H]
  \centering\scriptsize
  \setlength{\tabcolsep}{4pt}
  \caption{Đối chuẩn trên AMZN: 8 mô hình và mốc ngây thơ}
  \label{tab:bench-amzn}
  \input{generated/tab_bench_amzn}
\end{table}

\begin{table}[H]
  \centering\scriptsize
  \setlength{\tabcolsep}{4pt}
  \caption{Đối chuẩn trên KKBox: 8 mô hình}
  \label{tab:bench-kkbox}
  \input{generated/tab_bench_kkbox}
\end{table}

\begin{figure}[H]
  \centering
  \hinh{f4_scatter}
  \caption{Chất lượng theo độ trễ suy luận của 16 mô hình}
  \label{fig:scatter}
\end{figure}

Bảng~\ref{tab:bench-amzn}, Bảng~\ref{tab:bench-kkbox} và Hình~\ref{fig:scatter} cho bốn kết luận.

\textbf{Chất lượng gần như không phụ thuộc framework.} Trên KKBox, F1 của cùng một kiến trúc ở hai framework chênh nhau
từ \SL{bench/cmp/kkbox_rnn/diffPct} (Simple RNN) đến \SL{bench/cmp/kkbox_gru/diffPct} (GRU) điểm phần trăm. Trên AMZN, RMSE chênh không
quá \SL{bench/cmp/amzn_bilstm/diff} USD, trừ Simple RNN của Keras kém hơn \SL{bench/cmp/amzn_rnn/diff} USD vì dừng ở
epoch \SL{bench/models/amzn_keras_rnn/best_epoch}. Mỗi mô hình chỉ chạy một seed, nên những chênh lệch cỡ này nằm trong mức
dao động của khởi tạo ngẫu nhiên, không đủ để nói framework nào học tốt hơn.

\textbf{PyTorch nhanh hơn với từng mẫu một.} Độ trễ trung vị của 8 mô hình PyTorch là \SL{bench/lat_median/torch} ms,
của Keras là \SL{bench/lat_median/keras} ms. Với lô một mẫu, chi phí gọi hàm của framework chiếm phần lớn thời gian, và
lời gọi eager của PyTorch nhẹ hơn lời gọi \texttt{tf.function}. BiLSTM chênh nhiều nhất: Keras chậm hơn
\SL{bench/cmp/amzn_bilstm/lat_ratio} lần.

\textbf{Với lô lớn, khoảng cách hẹp lại và có chỗ đảo chiều.} Ở lô 256, Keras Simple RNN xử lý
\SL{bench/models/kkbox_keras_rnn/throughput} mẫu/giây trên KKBox, nhiều hơn PyTorch
(\SL{bench/models/kkbox_torch_rnn/throughput}); Keras GRU trên KKBox cũng nhanh hơn
(\SL{bench/models/kkbox_keras_gru/throughput} so với \SL{bench/models/kkbox_torch_gru/throughput}). Với LSTM và BiLSTM,
PyTorch vẫn nhanh hơn, từ \SL{bench/cmp/kkbox_bilstm/thr_ratio} đến \SL{bench/cmp/kkbox_lstm/thr_ratio} lần. Chọn framework theo tốc độ cần biết trước ứng dụng phục vụ từng
yêu cầu một hay theo lô.

\textbf{Mô hình ``tốt nhất'' khác nhau tuỳ tiêu chí.} Theo val loss (tiêu chí dùng để chọn mô hình cho Web App), mô hình
tốt nhất trên AMZN là \SL{bench/best/amzn_torchLabel} (PyTorch) và \SL{bench/best/amzn_kerasLabel} (Keras), trên
KKBox là \SL{bench/best/kkbox_torchLabel} (PyTorch) và \SL{bench/best/kkbox_kerasLabel} (Keras). Chỉ số test tốt nhất
lại thuộc về \SL{bench/best_test/amznLabel} trên AMZN và \SL{bench/best_test/kkboxLabel} trên KKBox. Khi các mô hình gần bằng nhau như vậy, thứ hạng trên test gần
như là ngẫu nhiên, và mô hình rẻ nhất đủ tốt là lựa chọn hợp lý.

\section{Giới hạn của mô hình dự báo giá}

Trên AMZN, không mô hình nào trong 16 mô hình thắng mốc ngây thơ (dòng cuối Bảng~\ref{tab:bench-amzn}). Đây không phải
lỗi cài đặt. Nó đến từ ba điều đã thấy trong dữ liệu:

\begin{itemize}
  \item \textbf{Return gần như không có tự tương quan} (Chương~\ref{ch:dulieu}): quá khứ 30 phiên chứa rất ít thông tin
        về return ngày mai. Dự báo tốt nhất theo sai số bình phương khi đó là một hằng số gần 0, đúng là điều mô hình
        học được (val loss đứng yên từ epoch đầu).
  \item \textbf{Trễ pha tại điểm đảo chiều}: vì $\hat r \approx 0$, dự báo $\hat P_{t+1} \approx P_t$ luôn đi sau giá thật
        một phiên (Hình~\ref{fig:pred-amzn}b). Khi giá đổi chiều, mô hình chỉ ``thấy'' sau khi việc đã xảy ra.
  \item \textbf{Phân phối đổi theo thời gian}: độ lệch chuẩn của return là \SL{amzn/ret_std/train} trên train nhưng chỉ
        \SL{amzn/ret_std/test} trên test (Chương~\ref{ch:dulieu}), và những phiên lớn nhất đến từ sự kiện ngoài dữ liệu (báo cáo tài chính, tin vĩ mô) mà
        tám đặc trưng kỹ thuật không chứa.
\end{itemize}

Chỉ số $R^2$ trên giá (khoảng 0,99) và biểu đồ ``đường dự báo bám sát đường thật'' là hai cách trình bày dễ làm kết quả
trông tốt hơn thực tế. Với chuỗi giá, cần so với mốc ngây thơ và đo trên return.

% =====================================================================
%  Chương 5 — Hạng mục 5: đóng gói triển khai Web App
%  Mã nguồn: app/app.py, app/static/app.js, app/smoke_test.py
% =====================================================================
\chapter{Triển khai Web App}
\label{ch:webapp}

Chương này đóng gói các mô hình thành một Web App chạy cục bộ. Ứng dụng có hai phân hệ ứng với hai tập dữ liệu: dự báo
giá đóng cửa AMZN của phiên kế tiếp, và dự báo xác suất churn của một người dùng KKBox từ chuỗi 31 ngày nghe nhạc. Người
dùng chọn framework (PyTorch hoặc Keras) và một trong bốn mô hình; mặc định là mô hình có val loss nhỏ nhất
(Chương~\ref{ch:keras}).

\section{Kiến trúc}

Máy chủ là Flask, trả JSON qua bốn nhóm endpoint; giao diện là một trang HTML với hai tab, gọi API bằng \texttt{fetch}
và vẽ biểu đồ bằng Chart.js (Hình~\ref{fig:arch}). Ứng dụng không cần dữ liệu gốc của KKBox để chạy. Mọi thứ nó cần
nằm trong thư mục \texttt{models/}: 16 tệp mô hình và \texttt{metadata.json} do notebook 06 ghi, gồm thống kê chuẩn hoá
của hai tập, ngưỡng phân loại chọn trên validation, tên mô hình tốt nhất và chuỗi hành vi của hai người dùng mẫu.

\begin{figure}[H]
  \centering
  \hinh{f5_arch}
  \caption{Kiến trúc Web App}
  \label{fig:arch}
\end{figure}

Bảng~\ref{tab:api} liệt kê các endpoint.

\begin{table}[H]
  \centering\small
  \caption{REST API của Web App}
  \label{tab:api}
  \begin{tabular}{llL{7.2cm}}
    \toprule
    Phương thức & Đường dẫn & Vào $\to$ ra \\
    \midrule
    GET & \texttt{/api/stock/history} & $\to$ 100 phiên gần nhất: ngày, Close, MA20 \\
    POST & \texttt{/api/stock/predict} & \texttt{\{framework, model\}} $\to$ giá dự báo, \% thay đổi, độ trễ \\
    POST & \texttt{/api/churn/predict} & \texttt{\{preset\}} hoặc \texttt{\{sequence\}} ($31\times8$), \texttt{framework},
                                        \texttt{model} $\to$ xác suất, mức rủi ro, ngưỡng, độ trễ \\
    GET & \texttt{/api/meta} & $\to$ danh sách mô hình, mô hình tốt nhất, người dùng mẫu \\
    GET & \texttt{/} & $\to$ trang giao diện \\
    \bottomrule
  \end{tabular}
\end{table}

\section{Nạp mô hình}

Mỗi mô hình chỉ được nạp một lần, ở lần gọi đầu tiên, rồi giữ trong bộ nhớ nhờ \texttt{lru\_cache}. Hàm nạp trả về một
hàm dự báo cùng giao diện cho cả hai framework, nên phần xử lý yêu cầu không cần biết mô hình đến từ đâu. Mô hình Keras
chạy qua \texttt{tf.function} như lúc đo độ trễ ở Chương~\ref{ch:keras}, và mỗi mô hình được chạy thử một lần ngay khi
nạp để độ trễ trả về cho người dùng là độ trễ lúc ổn định. TensorFlow chỉ được import khi có yêu cầu dùng mô hình Keras.

\maNguon{app__load_model}{Nạp mô hình một lần, cùng giao diện cho PyTorch và Keras}

\section{Hai phân hệ dự báo}

Phân hệ giá lấy 30 phiên mới nhất trong \texttt{AMZN.csv}, áp đúng phép biến đổi của Chương~\ref{ch:dulieu} (giá tương
đối, log Volume, z-score bằng thống kê của train), đưa qua mô hình rồi quy đổi return dự báo ra giá USD. Hàm
\texttt{latest\_window} dùng chung hai hàm biến đổi với \texttt{make\_windows} của lúc huấn luyện, nên dữ liệu phục vụ
không thể lệch khỏi dữ liệu huấn luyện.

\maNguon{app__stock_predict}{Endpoint dự báo giá phiên kế tiếp}

Phân hệ churn nhận hoặc tên một người dùng mẫu, hoặc một ma trận $31 \times 8$ bất kỳ. Đầu vào sai shape bị từ chối với mã
400 và thông báo lỗi rõ ràng, không đưa vào mô hình. Xác suất được so với ngưỡng của đúng mô hình đó để xếp mức rủi ro:
cao khi vượt ngưỡng, trung bình khi vượt nửa ngưỡng, còn lại là thấp.

\maNguon{app__churn_predict}{Endpoint dự báo churn}

\section{Giao diện}

Mỗi tab có một cột điều khiển và một biểu đồ. Khi đổi framework, danh sách mô hình được dựng lại và mô hình tốt nhất
trên validation của framework đó được chọn sẵn.

\maNguonJS{initChurn}{Tab churn: gọi API và vẽ chuỗi hành vi bằng Chart.js}

Hình~\ref{fig:ui-stock} và~\ref{fig:ui-churn} vẽ lại hai tab bằng TikZ thay cho ảnh chụp màn hình. Mọi con số và đường
cong trong hai hình là phản hồi thật của API, do \texttt{smoke\_test.py} ghi lại: biểu đồ giá đọc 100 phiên mà
\texttt{/api/stock/history} trả về, còn biểu đồ churn đọc đúng chuỗi mà \texttt{/api/churn/predict} trả về cho người dùng
mẫu.

\begin{figure}[H]
  \centering
  \hinh{f5_ui_stock}
  \caption{Tab dự báo giá AMZN với mô hình PyTorch tốt nhất trên validation}
  \label{fig:ui-stock}
\end{figure}

Với dữ liệu kết thúc ngày \SL{app/stock/torch/last_date}, mô hình \SL{app/stock/torch/model_label} của PyTorch dự báo
phiên kế tiếp đóng cửa ở \SL{app/stock/torch/pred_close} USD, mô hình \SL{app/stock/keras/model_label} của Keras dự báo
\SL{app/stock/keras/pred_close} USD, so với giá hiện tại \SL{app/stock/torch/last_close} USD. Cả hai chỉ lệch vài phần
nghìn so với giá hiện tại, đúng với điều Chương~\ref{ch:keras} đã phân tích: mô hình dự báo return gần 0.

\begin{figure}[H]
  \centering
  \hinh{f5_ui_churn}
  \caption{Tab dự báo churn với một người dùng thật đã churn trong tập test}
  \label{fig:ui-churn}
\end{figure}

Hai người dùng mẫu là người thật trong tập test: một người đã churn mà mô hình tin chắc nhất, và một người ở lại mà mô hình
cho xác suất churn thấp nhất. Người churn trong Hình~\ref{fig:ui-churn} nghe gần như mỗi ngày trong 9 ngày đầu tháng rồi
ngừng hẳn từ ngày 10, dạng cực đoan của quỹ đạo nhóm churn ở Chương~\ref{ch:dulieu}. LSTM của PyTorch cho xác suất
\SL{app/churn/churn_torch/probPct}\,\% với người này và \SL{app/churn/loyal_torch/probPct}\,\% với người ở lại; GRU của
Keras cho lần lượt \SL{app/churn/churn_keras/probPct}\,\% và \SL{app/churn/loyal_keras/probPct}\,\%.

\section{Kiểm chứng}

\texttt{smoke\_test.py} dùng test client của Flask, không cần mở cổng mạng. Nó gọi lịch sử giá, dự báo giá với cả hai
framework và dự báo churn cho hai người dùng mẫu với cả hai framework, tổng cộng \SL{app/n_calls} lượt. Mỗi phản hồi phải
có mã 200 và đủ các khoá cần thiết; một chuỗi sai shape phải bị từ chối với mã \SL{app/bad_shape_status}. Kịch bản dừng
ngay khi có một điều kiện không thoả, và lượt chạy cuối qua hết.

\maNguon{smoke_test__main}{Kiểm chứng ba endpoint và ghi phản hồi thật cho báo cáo}

Chạy ứng dụng bằng \texttt{python app/app.py} rồi mở \texttt{http://localhost:8080}. Độ trễ API trả về là thời gian của riêng lời gọi mô hình:
\SL{app/churn/churn_torch/latency_ms} ms với LSTM của PyTorch và \SL{app/churn/churn_keras/latency_ms} ms với GRU của
Keras, cùng cỡ với số đo ở Bảng~\ref{tab:bench-kkbox}. Trong ứng dụng, PyTorch chạy 1 luồng còn TensorFlow dùng số luồng
mặc định, và mỗi con số chỉ là một lần gọi, nên chúng để tham khảo chứ không thay cho phép đo có kiểm soát của
Chương~\ref{ch:keras}.

% =====================================================================
%  Chương 6 — Kết luận
% =====================================================================
\chapter{Kết luận}
\label{ch:ketluan}

\section{Trả lời sáu hạng mục}

\begin{enumerate}
  \item \textbf{Cơ sở lý thuyết.} Mỗi khái niệm của mạng hồi quy (tensor chuỗi, Simple RNN, BPTT, clipping, tế bào LSTM và
        GRU, BiLSTM, recurrent dropout, LayerNorm) được viết thành một hàm NumPy và khớp với PyTorch ở cả
        \SL{theory/n_checks} phép đối chiếu. Phép đo gradient bằng autograd cho một kết quả đáng nhớ: LSTM chỉ giữ được
        gradient khi cổng quên mở. Với bias cổng quên bằng 0, gradient ở trễ 20 gần như bằng của Simple RNN.
  \item \textbf{Hai tập dữ liệu thực tế.} AMZN có \SL{amzn/n_raw} phiên, chia theo thời gian, dự báo return thay vì mức
        giá vì giá không dừng (ADF $p = \SL{amzn/adf/close/p}$). KKBox có \SL{kkbox/n_sample} người dùng thật lấy từ
        \SL{kkbox/log_rows} dòng log, đọc theo khối không cần nạp trọn vào RAM. Trên cả hai tập, mọi thống kê chuẩn hoá
        chỉ tính trên train.
  \item \textbf{PyTorch.} Bốn mô hình dùng chung một lớp và một vòng lặp huấn luyện, huấn luyện trên GPU sau khi đo được
        GPU nhanh hơn CPU. Trên KKBox, LSTM tốt nhất với F1 \SL{torch_kkbox/models/lstm/f1}\,\% và ROC-AUC
        \SL{torch_kkbox/models/lstm/roc_auc}; ba mô hình có cổng đều hơn Simple RNN.
  \item \textbf{Keras và đối chuẩn.} Cùng kiến trúc cho chất lượng gần như bằng nhau ở hai framework. PyTorch nhanh hơn
        với từng mẫu một (trung vị \SL{bench/lat_median/torch} so với \SL{bench/lat_median/keras} ms trên một luồng CPU),
        còn với lô lớn thì khoảng cách hẹp lại và Keras thắng ở Simple RNN. Tệp \texttt{.keras} nặng hơn vì chứa cả trạng thái
        optimizer.
  \item \textbf{Web App.} Flask phục vụ cả 16 mô hình qua REST API, giao diện hai tab; toàn bộ tham số tiền xử lý đọc từ
        \texttt{metadata.json} nên ứng dụng chạy được mà không cần dữ liệu gốc. Kịch bản kiểm chứng qua
        \SL{app/n_calls} lượt gọi.
  \item \textbf{Kết luận về mô hình.} Mạng hồi quy học được tín hiệu thật khi tín hiệu nằm trong chuỗi (xu hướng nghe nhạc
        giảm dần của người sắp rời bỏ). Nó không tạo ra được tín hiệu khi chuỗi không chứa (return cổ phiếu ngày mai):
        không mô hình nào trong 16 mô hình thắng mốc ``giá ngày mai bằng giá hôm nay''.
\end{enumerate}

\section{Hạn chế}

\begin{itemize}
  \item Mỗi mô hình chỉ chạy một seed. Chênh lệch giữa các mô hình trên KKBox (dưới hai điểm F1) và giữa hai framework
        đều nhỏ hơn mức dao động thường gặp giữa các lần khởi tạo, nên bài không xếp hạng chúng.
  \item KKBox chỉ dùng mẫu 100.000 người và 31 ngày log. Bỏ qua lịch sử giao dịch và thông tin thành viên nên ROC-AUC dừng
        ở khoảng 0,68, thấp hơn các lời giải đầy đủ của cuộc thi.
  \item AMZN chỉ dùng đặc trưng kỹ thuật tính từ chính chuỗi giá. Các sự kiện gây biến động lớn nhất (báo cáo quý, tin vĩ
        mô) không có trong dữ liệu.
  \item Keras chạy trên CPU vì TensorFlow không hỗ trợ GPU trên Windows, nên thời gian huấn luyện của hai framework không
        so trực tiếp được; bài chỉ so tốc độ suy luận trên cùng CPU.
\end{itemize}


% =====================================================================
%  TÀI LIỆU THAM KHẢO, PHỤ LỤC
% =====================================================================
\cleardoublepage
\phantomsection
\addcontentsline{toc}{chapter}{Tài liệu tham khảo}
{\sloppy\printbibliography[title={Tài liệu tham khảo}]}

\appendix
% =====================================================================
%  Phụ lục — Tái lập kết quả
% =====================================================================
\chapter{Tái lập kết quả}
\label{ch:phuluc}

\section{Môi trường}

Windows 11, Python 3.13, GPU NVIDIA GeForce RTX 4060 Laptop cho PyTorch (CUDA 12.6), CPU cho TensorFlow. PyTorch
\SL{bench/torch_version}, Keras \SL{bench/keras_version} với TensorFlow 2.21. Mọi thư viện cài trong môi trường ảo riêng
của bài, liệt kê trong \texttt{requirements.txt}.

\section{Các bước}

\begin{enumerate}
  \item Tạo môi trường và cài thư viện:
\begin{verbatim}
cd "src/Assignment 06"
python -m venv .venv
.venv\Scripts\activate
pip install torch==2.13.0 --index-url https://download.pytorch.org/whl/cu126
pip install -r requirements.txt
python -m ipykernel install --user --name a06
\end{verbatim}
  \item Tải \texttt{train\_v2.csv.7z} và \texttt{user\_logs\_v2.csv.7z} của cuộc thi KKBox trên Kaggle (cần đăng nhập và
        chấp nhận luật cuộc thi) vào \texttt{data/kkbox/}, giải nén. \texttt{AMZN.csv} đã có sẵn trong kho.
  \item Chạy các notebook theo thứ tự số. Notebook 03, 04 (GPU) và 05 (CPU) chạy song song được. Trước khi chạy, đặt
        biến môi trường \texttt{OPENBLAS\_NUM\_THREADS} bằng 8: trên máy 32 luồng, OpenBLAS mặc định cấp bộ đệm cho 32
        luồng và có lúc làm kernel chết vì không cấp phát được bộ nhớ.
\begin{verbatim}
cd notebooks
jupyter nbconvert --to notebook --execute --inplace ^
    --ExecutePreprocessor.kernel_name=a06 --ExecutePreprocessor.timeout=-1 00_*.ipynb
\end{verbatim}
  \item Kiểm chứng Web App: \texttt{python app/smoke\_test.py}; chạy thật: \texttt{python app/app.py}.
  \item Dựng báo cáo: \texttt{python "Report/Assignment 06/build.py"} (cần LuaLaTeX, \texttt{latexmk}, \texttt{biber}
        và gói \texttt{minted}).
\end{enumerate}

Lần chạy đầu của notebook 02 gom log KKBox thành tệp \texttt{kkbox\_seq.npz}; các lần sau đọc thẳng tệp này.

\maNguon{data_churn__build_cache}{Gom dữ liệu KKBox một lần và lưu lại}


\end{document}
